\notepage{9}{The eigenvalue problem}
In general, $\Omega\ket{v}=\ket{v'}$ changes the direction of a vector.
An \emph{eigenvector} is a nonzero vector whose direction is unchanged:
\[
 \Omega\ket{v}=\omega\ket{v},\qquad \ket{v}\ne0.
\]
$\omega$ is its eigenvalue. For example, $I\ket{v}=\ket{v}$ for every
vector, so every nonzero vector is an eigenvector of $I$ with eigenvalue~1.
If $\Omega\ket{v}=\omega\ket{v}$, linearity gives
\[
 \Omega(\alpha\ket{v})=\omega(\alpha\ket{v})
\]
for every nonzero scalar $\alpha$.

\smallheading{The characteristic equation}
The eigenvalue equation is equivalent to
\[
 (\Omega-\omega I)\ket{v}=0.
\]
If $\Omega-\omega I$ had an inverse, multiplication by that inverse
would give $\ket{v}=0$, which is excluded. Thus, in finite dimensions,
\[
 \boxed{\det(\Omega-\omega I)=0.}
\]
To write this in a basis, insert the identity:
\begin{align*}
 0&=\mel{i}{\Omega-\omega I}{v}\\
  &=\sum_j\mel{i}{\Omega-\omega I}{j}\braket{j}{v}\\
  &=\sum_j\bigl(\Omega_{ij}-\omega\delta_{ij}\bigr)v_j,
 \qquad v_j=\braket{j}{v}.
\end{align*}
The determinant condition becomes
\[
 \begin{vmatrix}
 \Omega_{11}-\omega&\Omega_{12}&\cdots&\Omega_{1n}\\
 \Omega_{21}&\Omega_{22}-\omega&\cdots&\Omega_{2n}\\
 \vdots&\vdots&\ddots&\vdots\\
 \Omega_{n1}&\Omega_{n2}&\cdots&\Omega_{nn}-\omega
 \end{vmatrix}
 =\sum_{m=0}^{n}c_m\omega^m=0.
\]
This polynomial has degree $n$. Its roots are the eigenvalues, which
are independent of the basis used to calculate the determinant.
